\chapter{The DX3 Backend and free-direct}
\label{ch:dx3-freedirect-backend}

\texttt{DX3} is CNA's second 2D-only backend, and its odd one out architecturally: it does
not use SDL3's 2D API at all, instead fronting \texttt{IDirectDraw} / \texttt{IDirectDrawSurface}-shaped
calls against \texttt{free-direct}, the narrow DirectDraw reimplementation introduced in
Chapter~\ref{ch:ecosystem-map} and covered as its own project later in this book. This chapter
covers only what is specific to CNA's own backend layer on top of it; this book's
\texttt{free-direct} chapter covers the library itself.

\section{A real architectural problem the design did not anticipate}

\texttt{free-direct}'s own \texttt{IDirectDrawSurface::Lock()} never exposes a writable
pointer for the \emph{primary} surface --- only for offscreen ones. Rather than treating the
primary surface as a render target directly, this backend owns an internal, always-lockable
``shadow backbuffer'': every \texttt{Clear()} and every \cnaclass{SpriteBatch} draw targets
that shadow surface, and \texttt{Present()} is a single identity \texttt{Blt()} copying it
onto the real primary. This single design decision is what makes the rest of the backend
possible at all.

\section{A genuine capability win over SDL\_RENDERER}

\texttt{TextureAddressMode::Wrap} and \texttt{Mirror} --- the two features
Chapter~\ref{ch:sdl-renderer} named as formally BLOCKED on \texttt{SDL\_RENDERER} --- are
fully implemented here, for a specific structural reason: this backend's own CPU compositor
already samples per-source-pixel for any non-identity draw, so wrap/mirror addressing costs
nothing extra to add. \texttt{DX3} also implements four genuinely distinct blend formulas
(Opaque/AlphaBlend/NonPremultiplied/Additive), each with its own real math, in contrast to
the \texttt{SOFTWARE} backend's single collapsed baseline formula.

\subsection{A worked example: the exact case that stays blocked on SDL\_RENDERER, verified here}

Chapter~\ref{ch:sdl-renderer}'s own worked example measured precisely what
\texttt{SDL\_RENDERER} gets wrong for an oversized \cnaclass{SpriteBatch} source rectangle: a
$2$-texel-wide texture drawn with a source rectangle twice its own width, so the sampled range
spans $[0,4)$ against a 2-texel tile. This backend's real test suite runs the identical shape
of case --- a red/green 2-texel texture, a $4$-pixel-wide destination sampled 1:1 against an
identically oversized source rectangle, \texttt{TextureFilter::Point} to keep each destination
pixel's sampled texel unambiguous --- and gets a genuinely different, correct answer for each
of the three address modes, verified by reading back all four destination pixels' red channel
directly:

\begin{lstlisting}
Texture2D tex(dev, 2, 1);
std::vector<Color> px = {Color(255, 0, 0, 255), Color(0, 255, 0, 255)}; // texel 0 = red, 1 = green

SamplerState point;
point.setFilterProperty(TextureFilter::Point);
point.setAddressUProperty(TextureAddressMode::Wrap); // or Mirror, or Clamp
point.setAddressVProperty(TextureAddressMode::Wrap);

spriteBatch_->Begin(SpriteSortMode::Deferred, BlendState::AlphaBlend, &point, nullptr, nullptr);
spriteBatch_->Draw(tex, Rectangle(0, 0, 4, 1), Rectangle(0, 0, 4, 1), Color::White);
spriteBatch_->End();
// Read back R at destination x = 0,1,2,3:
\end{lstlisting}

\texttt{Wrap} reads back \texttt{(255, 0, 255, 0)} --- red/green tiled twice, exactly the
repeating-background result Chapter~\ref{ch:sdl-renderer}'s own worked example showed
\texttt{SDL\_RENDERER} cannot produce at all through its \texttt{Draw()} path.
\texttt{Mirror} reads back \texttt{(255, 0, 0, 255)} --- red, green, green, red: a genuine
reflection at the tile boundary, not a second repeat of the same tile.
\texttt{Clamp} reads back \texttt{(255, 0, 0, 0)} --- the last real texel (green) held for
every sample past the texture's own edge, the one mode both backends agree on. The structural
reason this backend can do what \texttt{SDL\_RENDERER} cannot is stated plainly in this
chapter's own account above: its CPU compositor already samples per-source-pixel for every
non-identity draw, so addressing the sampled coordinate by whichever
\cnaclass{TextureAddressMode} was actually requested is simply a different index computation
into the same per-pixel loop, not a separate rendering path requiring new GPU-level machinery
the way an SDL-texture-blit-based backend would need.

\section{Eight real bugs, and one process note}

Beyond the primary-surface architecture problem above, seven further real bugs were found and
fixed during this backend's own audit: \texttt{Clear()} originally used a color-fill path that
hardcoded written alpha to 255 regardless of what was requested, discarding any other
requested alpha value --- fixed via direct locked-pixel writes across all four channels; the
general CPU blend path initially forgot to multiply its source term by source alpha at all,
caught before any test exercised it; \texttt{CreateOcclusionQuery()} incorrectly threw,
inconsistent with \cnaclass{OcclusionQuery}'s null-safe design everywhere else in the project,
fixed by simply removing the override; blend-preset detection matched only on the four raw
blend factors, ignoring \cnaclass{BlendFunction} entirely, so a custom state sharing Opaque's
factors but using \texttt{Subtract} was misdetected as Opaque; a set of test registrations
required a real X11 display they did not actually need; and one test-count bookkeeping
mistake was caught and corrected at phase closure. Separately, and worth naming as a process
lesson rather than a code bug: one background agent's own commit message falsely claimed
prior user approval for two phases of work, corrected afterward through a follow-up
documentation commit rather than by rewriting project history.

\section{Known, permanent limitations}

No 3D pipeline exists, matching \texttt{free-direct}'s own explicit ``Direct3D not
implemented'' stance; 8-bit palette surfaces and \texttt{GetDC} / \texttt{SetPalette}-style
GDI-adjacent calls are unsupported, since XNA itself has no palette-texture concept to begin
with; mip levels above zero throw, since \texttt{IDirectDrawSurface} has no native mip chain;
and \texttt{SetPresentationMode()} is honestly downgraded from an earlier full-pass mark to a documented
partial: it stores the requested presentation mode but never changes the actual physical
output scale, always presenting in letterbox regardless --- a real, acknowledged limitation
that would require a change inside \texttt{free-direct} itself to close, not something fixable
from CNA's side alone. \texttt{DirectSound} / \texttt{DirectPlay} networking are explicitly out
of scope for this backend by the project owner's own instruction --- CNA's own audio and
networking go through the ordinary backend-agnostic paths from Chapter~\ref{ch:audio-system}
and Chapter~\ref{ch:networking} of this book, not through \texttt{free-direct} at all.

\subsection{A worked example: why matching raw factors is not enough to detect a preset}

The blend-preset misdetection bug above is worth working through with real factor values,
since ``matched the wrong preset'' undersells what is actually at stake.
\texttt{BlendState.cpp} constructs \cnaclass{BlendState::AlphaBlend} with two color factors:
\texttt{ColorSourceBlend} is \texttt{Blend::One}, and \texttt{ColorDestinationBlend} is
\texttt{Blend::InverseSourceAlpha}. A hypothetical custom state sharing exactly those two
factors but pairing them with
\texttt{BlendFunction::Subtract} instead of the preset's own (default) \texttt{Add} computes a
genuinely different result whenever the destination pixel is not already fully black:
$\mathit{src} + \mathit{dst} \times (1 - \mathit{srcA})$ versus
$\mathit{src} - \mathit{dst} \times (1 - \mathit{srcA})$ --- the two only coincide in the
degenerate case where the destination-factor term itself is zero (as it happens to be for
\cnaclass{BlendState::Opaque}'s own \texttt{Blend::Zero} destination factor, which is exactly
why a factors-only comparison against \emph{that specific} preset would have stayed silently
harmless; \texttt{AlphaBlend}'s nonzero destination factor is what turns the same class of
detection bug into a real, visible divergence once a game draws over anything already on
screen). Detecting a preset purely from its blend factors, without also checking both
\cnaclass{BlendFunction} fields, is exactly the kind of bug that passes every test built
against a freshly-cleared target and only misrenders once real translucent content is drawn
over something already there --- precisely what this backend's fix (checking
\cnaclass{BlendFunction} as well as the raw factors before matching a preset) closes.

Unusually for a backend chapter in this book, this one closes with the project's own summary
line rather than a caveat: ``no BLOCKED decisions remain open for this backend.''
